% Compilation : xelatex RAPPORT.tex   (deux fois pour la table des matieres)
\documentclass[11pt,a4paper]{article}

\usepackage{fontspec}
\usepackage{polyglossia}
\setdefaultlanguage{french}
\usepackage{geometry}
\geometry{margin=2.3cm}
\usepackage{xcolor}
\usepackage{listings}
\usepackage{booktabs}
\usepackage{array}
\usepackage{graphicx}
\usepackage{float}
\usepackage{enumitem}
\usepackage{fancyhdr}
\usepackage[hidelinks]{hyperref}

% ---------- Couleurs ----------
\definecolor{codebg}{RGB}{245,245,245}
\definecolor{codeframe}{RGB}{210,210,210}
\definecolor{kw}{RGB}{0,0,160}
\definecolor{cmt}{RGB}{0,128,0}
\definecolor{str}{RGB}{163,21,21}
\definecolor{insa}{RGB}{200,16,46}

% ---------- Style du code Java ----------
\lstdefinestyle{java}{
  language=Java,
  backgroundcolor=\color{codebg},
  basicstyle=\ttfamily\footnotesize,
  keywordstyle=\color{kw}\bfseries,
  commentstyle=\color{cmt}\itshape,
  stringstyle=\color{str},
  numbers=left,
  numberstyle=\tiny\color{gray},
  numbersep=8pt,
  frame=single,
  rulecolor=\color{codeframe},
  breaklines=true,
  showstringspaces=false,
  tabsize=4,
  extendedchars=true,
  columns=fullflexible,
  keepspaces=true,
  literate={é}{{\'e}}1 {è}{{\`e}}1 {ê}{{\^e}}1 {à}{{\`a}}1 {â}{{\^a}}1
           {î}{{\^i}}1 {ô}{{\^o}}1 {û}{{\^u}}1 {ù}{{\`u}}1 {ç}{{\c c}}1
           {É}{{\'E}}1 {È}{{\`E}}1 {Ê}{{\^E}}1 {À}{{\`A}}1 {«}{{«}}1 {»}{{»}}1
}
\lstset{style=java}

% ---------- En-tete / pied ----------
\pagestyle{fancy}
\fancyhf{}
\lhead{\small TD 3 — Programmation concurrente en Java}
\rhead{\small STI 4A-FISE 2026/2027}
\cfoot{\thepage}
\renewcommand{\headrulewidth}{0.4pt}

% ---------- Macro captures d'ecran ----------
% Insere captures/<fichier> s'il existe, sinon un cadre "a inserer".
\newcommand{\capture}[2]{%
  \begin{figure}[H]\centering
  \IfFileExists{captures/#1}{%
    \includegraphics[width=0.85\textwidth]{captures/#1}%
  }{%
    \fbox{\parbox[c][3.2cm][c]{0.85\textwidth}{\centering\itshape
      \color{gray} Capture d'écran à insérer : \texttt{captures/#1}}}%
  }%
  \caption{#2}
  \end{figure}%
}

\begin{document}

% ================= PAGE DE TITRE =================
\begin{titlepage}
\centering
\vspace*{1cm}
\includegraphics[width=0.5\textwidth]{assets/logo-insa.png}\\[2cm]
{\color{insa}\rule{\textwidth}{2pt}}\\[0.6cm]
{\Huge\bfseries TD 3\par}
\vspace{0.4cm}
{\LARGE Programmation Concurrente en Java\par}
\vspace{0.3cm}
{\large Programmation Java Avancée : POO et JVM\par}
{\color{insa}\rule{\textwidth}{2pt}}\\[1.5cm]

{\large\itshape Rapport de travaux pratiques\par}
\vspace{1cm}

{\Large\bfseries Ali Jamai\par}
\vspace{1.5cm}

\begin{tabular}{rl}
\textbf{Auteur} & Ali Jamai\\[4pt]
\textbf{Module} & STI 4A-FISE 2026/2027\\[4pt]
\textbf{Enseignante} & Dr. Nadia Badri Fettah\\[4pt]
\textbf{Durée} & 1h20\\[4pt]
\textbf{Code source} & \url{https://github.com/JamaiAli/td3java}\\
\end{tabular}

\vfill
{\large \today\par}
\end{titlepage}

\tableofcontents
\newpage

% ================= OBJECTIFS =================
\section*{Objectifs du TD}
\addcontentsline{toc}{section}{Objectifs du TD}
\begin{itemize}[nosep]
  \item Créer et démarrer des threads avec \texttt{Runnable} et avec \texttt{extends Thread}.
  \item Observer le cycle de vie d'un thread (\texttt{getState}, \texttt{sleep}, \texttt{join}, \texttt{daemon}).
  \item Mettre en évidence une \emph{race condition} et la corriger avec \texttt{synchronized}.
  \item Utiliser \texttt{volatile}, \texttt{AtomicInteger}, \texttt{ExecutorService}, \texttt{Callable} et \texttt{Future}.
  \item Synchroniser le départ de plusieurs threads avec \texttt{wait()} et \texttt{notifyAll()}.
\end{itemize}

\begin{quote}\small\itshape
Chaque capture d'écran provient d'une exécution réelle du programme dans IntelliJ
IDEA ; les valeurs citées dans le texte (durées, tailles, résultats) correspondent à
ces exécutions. Comme l'ordre d'exécution des threads varie, les valeurs exactes
peuvent différer d'une exécution à l'autre.
\end{quote}

% ================= EXERCICE 1 =================
\section{Téléchargements : séquentiel ou parallèle ?}
On simule 3 téléchargements : \texttt{photo.jpg} (2\,s), \texttt{video.mp4} (3\,s)
et \texttt{cours.pdf} (1\,s). Un « téléchargement » est une simple attente avec
\texttt{sleep()}.

\subsection{Question 1.1 — La tâche (\texttt{Telechargement.java})}
\begin{lstlisting}
public class Telechargement implements Runnable {
    private String fichier;
    private int secondes;

    public Telechargement(String fichier, int secondes) {
        this.fichier = fichier;
        this.secondes = secondes;
    }

    public void run() {
        String qui = Thread.currentThread().getName();
        System.out.println(qui + " commence " + fichier);
        for (int i = 1; i <= secondes; i++) {
            try {
                Thread.sleep(1000);
            } catch (InterruptedException e) { }
            System.out.println("   " + fichier + " : " + (i * 100 / secondes) + " %");
        }
        System.out.println(qui + " a terminé " + fichier);
    }
}
\end{lstlisting}

\subsection{Question 1.2 — Version séquentielle}
On appelle \texttt{run()} directement : tout est exécuté par le thread \texttt{main},
les fichiers se téléchargent l'un après l'autre. \textbf{Durée $\approx$ 6000\,ms
(2 + 3 + 1).}
\begin{lstlisting}
d1.run();
d2.run();
d3.run();
\end{lstlisting}

\subsection{Question 1.3 — Version parallèle (\texttt{Navigateur.java})}
On crée un thread par téléchargement, on les démarre avec \texttt{start()}, puis on
attend leur fin avec \texttt{join()}.
\begin{lstlisting}
Thread t1 = new Thread(d1, "DL-1");
Thread t2 = new Thread(d2, "DL-2");
Thread t3 = new Thread(d3, "DL-3");

t1.start(); t2.start(); t3.start();
t1.join();  t2.join();  t3.join();
\end{lstlisting}
À l'exécution (voir capture), les pourcentages des trois fichiers sont mélangés et
la durée totale est $\approx$ \textbf{3000\,ms} (le plus long des trois), contre
$\approx$ 6000\,ms en séquentiel.

\capture{ex1.png}{Exécution parallèle de \texttt{Navigateur} — durée $\approx$ 3000\,ms}

\subsection{Question 1.4 — Sans \texttt{join()}}
Si on met les trois \texttt{join()} en commentaire, la durée affichée tombe à
$\approx$ 0\,ms. C'est \textbf{faux} : le \texttt{main} n'attend plus la fin des
téléchargements, il calcule la durée immédiatement après avoir lancé les threads.

\subsection{Question 1.5 — Tableau comparatif}
\begin{table}[H]\centering
\begin{tabular}{lcc}
\toprule
\textbf{Version} & \textbf{Durée mesurée (ms)} & \textbf{Thread qui exécute} \\
\midrule
Séquentielle (\texttt{run()}) & $\approx$ 6000 & \texttt{main} \\
Parallèle (\texttt{start()}+\texttt{join()}) & $\approx$ 3008 & \texttt{DL-1}, \texttt{DL-2}, \texttt{DL-3} \\
\bottomrule
\end{tabular}
\end{table}
\textbf{Pourquoi la version parallèle est-elle plus rapide, même sur un seul cœur ?}
Parce que pendant un \texttt{sleep()}, un thread \textbf{n'utilise pas le
processeur}. Les trois attentes se déroulent en même temps au lieu de s'additionner,
d'où une durée égale au plus long ($\approx$ 3\,s) et non à leur somme ($\approx$ 6\,s).

% ================= EXERCICE 2 =================
\section{Cycle de vie : sleep, join, getState, daemon}

\subsection{Question 2.1 — Afficher les états (\texttt{Etats.java})}
\begin{lstlisting}
System.out.println("1. " + t.getState()); // avant start()  -> NEW
t.start();
System.out.println("2. " + t.getState()); // juste après start() -> RUNNABLE
Thread.sleep(200);
System.out.println("3. " + t.getState()); // t dort -> TIMED_WAITING
t.join();
System.out.println("4. " + t.getState()); // après fin -> TERMINATED
\end{lstlisting}
À l'exécution, les états s'enchaînent dans l'ordre attendu :
\texttt{NEW} $\rightarrow$ \texttt{RUNNABLE} $\rightarrow$ \texttt{TIMED\_WAITING}
$\rightarrow$ \texttt{TERMINATED} (voir capture).

\capture{ex2.png}{États successifs du thread (NEW, RUNNABLE, TIMED\_WAITING, TERMINATED)}

\subsection{Question 2.2 — Le rôle de \texttt{join()}}
Si on commente le \texttt{join()} juste avant la ligne 4, celle-ci peut afficher
\textbf{\texttt{TIMED\_WAITING}} (le thread dort encore) au lieu de
\texttt{TERMINATED}. \texttt{join()} \textbf{bloque le thread appelant jusqu'à ce que
le thread visé soit réellement terminé.}

\subsection{Question 2.3 — Redémarrer un thread terminé}
Ajouter \texttt{t.start();} à la fin du \texttt{main} lève
\texttt{java.lang.IllegalThreadStateException} : \textbf{un thread ne peut être
démarré qu'une seule fois}. Une fois \texttt{TERMINATED}, il faut créer un nouvel
objet \texttt{Thread}.

\subsection{Question 2.4 — Le thread daemon (\texttt{Horloge.java})}
\begin{itemize}[nosep]
  \item \textbf{Étape a} (sans \texttt{setDaemon}) : la boucle \texttt{while(true)}
  ne s'arrête jamais ; il faut \texttt{Ctrl + C}.
  \item \textbf{Étape b} (\texttt{t.setDaemon(true);} \emph{avant} \texttt{start()}) :
  $\sim$4 lignes \texttt{Tic}, puis \texttt{Fin du main}, puis arrêt automatique.
  Quand il ne reste que des threads daemon, la JVM se termine.
  \item \textbf{Étape c} (\texttt{setDaemon(true)} \emph{après} \texttt{start()}) :
  lève \texttt{IllegalThreadStateException}.
\end{itemize}

\capture{ex2b.png}{Horloge daemon : arrêt après « Fin du main »}

% ================= EXERCICE 3 =================
\section{Une ArrayList partagée qui perd des éléments}

\subsection{Question 3.1 — Le bug}
Deux threads ajoutent chacun 10\,000 entiers dans la même \texttt{ArrayList} non
protégée (taille attendue 20000). On obtient des tailles \textbf{inférieures à
20000} et parfois une \texttt{ArrayIndexOutOfBoundsException}.
\begin{table}[H]\centering
\begin{tabular}{ccc}
\toprule
\textbf{Exécution} & \textbf{Taille obtenue} & \textbf{Exception ?} \\
\midrule
1 (capture) & 12706 & non \\
2 & 11285 & non \\
3 & 10000 & oui (AIOOBE) \\
4 & 11442 & non \\
5 & 10004 & oui (AIOOBE) \\
\bottomrule
\end{tabular}
\end{table}
Les valeurs varient à chaque exécution (l'exécution 1 est celle de la capture). On
n'obtient jamais 20000, et une \texttt{ArrayIndexOutOfBoundsException} apparaît par
moments. \emph{Si vous obtenez toujours 20000, remplacez \texttt{10\_000} par
\texttt{100\_000} (taille attendue 200000).}

\capture{ex3.png}{Taille < 20000 et/ou ArrayIndexOutOfBoundsException}

\subsection{Question 3.2 — Comprendre}
\texttt{liste.add(x)} fait deux choses : (1) ranger \texttt{x} dans la case n°
\texttt{size} ; (2) \texttt{size++}. Ces étapes n'étant pas atomiques, deux threads
peuvent lire la \textbf{même} valeur de \texttt{size}, écrire dans la \textbf{même}
case (l'un écrase l'autre → élément perdu), puis incrémenter \texttt{size} deux fois
→ taille incohérente et case « vide » → exception.

\subsection{Question 3.3 — Correction n°1 : un bloc \texttt{synchronized}}
\begin{lstlisting}
synchronized (liste) {
    liste.add(i);
}
\end{lstlisting}
\textbf{Résultat :} \texttt{Taille obtenue : 20000} à chaque exécution, sans
exception (résultat attendu). On verrouille la \textbf{ressource partagée elle-même}
car le \texttt{main} est \texttt{static} (pas de \texttt{this}).

\subsection{Question 3.4 — Correction n°2 : une liste synchronisée}
\begin{lstlisting}
import java.util.Collections;
List<Integer> liste = Collections.synchronizedList(new ArrayList<>());
\end{lstlisting}
Chaque méthode (\texttt{add}, \texttt{size}, \texttt{get}…) est déjà
\texttt{synchronized}. Résultat : 20000 à chaque fois.

\subsection{Question 3.5 — Réflexion}
\begin{lstlisting}
if (!liste.contains(42)) {
    liste.add(42);
}
\end{lstlisting}
\textbf{Peut-on obtenir deux fois 42 ?} \textbf{Oui.} Chaque appel est protégé
individuellement, mais \textbf{pas l'ensemble « vérifier puis ajouter »}. Deux threads
peuvent évaluer \texttt{contains(42)} à \texttt{false} avant qu'aucun n'ait ajouté.
\textbf{Correction :} rendre le bloc composé atomique.
\begin{lstlisting}
synchronized (liste) {
    if (!liste.contains(42)) {
        liste.add(42);
    }
}
\end{lstlisting}

% ================= EXERCICE 4 =================
\section{Bloc \texttt{synchronized} et verrous distincts}

\subsection{Question 4.1 — Deux ressources, un seul verrou}
Avec \texttt{public synchronized void ...}, le verrou est \textbf{l'objet \texttt{b}
tout entier} : \texttt{ajouterLivre} et \texttt{ajouterLecteur} ne peuvent pas
s'exécuter en même temps. \textbf{Durée $\approx$ 2000\,ms} (20 appels × 100\,ms en
série). Les deux ressources sont indépendantes, pourtant \texttt{t2} attend
\texttt{t1} inutilement.

\subsection{Question 4.2 — Un verrou par ressource}
\begin{lstlisting}
public void ajouterLivre(String titre) {
    synchronized (livres) { livres.add(titre); }
    pause();                       // hors du bloc
}
public void ajouterLecteur(String nom) {
    synchronized (lecteurs) { lecteurs.add(nom); }
    pause();
}
\end{lstlisting}
À l'exécution, la durée tombe à \texttt{Durée : 1010 ms} $\approx$ \textbf{1000\,ms}
(voir capture), contre $\approx$ 2000\,ms avec un verrou unique : les deux threads
avancent maintenant en parallèle (verrous distincts).

\capture{ex4.png}{Durée $\approx$ 1000\,ms avec un verrou par ressource}

\subsection{Question 4.3 — Réflexion}
On laisse \texttt{pause()} hors de la section critique car garder un verrou pendant
une opération longue bloquerait inutilement les autres threads ; on ne verrouille que
la modification réellement partagée (l'ajout dans la liste).

% ================= EXERCICE 5 =================
\section{\texttt{volatile} et \texttt{AtomicInteger}}

\subsection{Question 5.1 — Arrêter proprement un thread (\texttt{Travailleur.java})}
Sans \texttt{volatile}, le thread peut \textbf{ne jamais voir} que \texttt{actif} est
passé à \texttt{false} (valeur mise en cache). Sur beaucoup de machines on lit alors
\texttt{Le thread tourne encore !}. Le problème n'est pas garanti mais possible — ce
qui le rend dangereux.

\subsection{Question 5.2 — La correction}
\begin{lstlisting}
private volatile boolean actif = true;
\end{lstlisting}
Toute écriture de \texttt{actif} devient immédiatement visible par les autres threads :
le thread s'arrête et le programme affiche \texttt{Le thread est bien arrêté} à chaque
exécution (voir capture).

\capture{ex5.png}{« Le thread est bien arrêté » avec \texttt{volatile}}

\subsection{Question 5.3 — \texttt{volatile} ne suffit pas pour \texttt{++}}
Avec \texttt{private volatile int vues = 0;} et \texttt{vues++}, le résultat est
\textbf{faux} (< 200000). \texttt{volatile} garantit la \textbf{visibilité} mais
\textbf{pas l'atomicité} : \texttt{vues++} = \emph{lire → +1 → écrire} (3 étapes).

\subsection{Question 5.4 — \texttt{AtomicInteger}}
\begin{lstlisting}
import java.util.concurrent.atomic.AtomicInteger;
private AtomicInteger vues = new AtomicInteger(0);
public void nouvelleVue() { vues.incrementAndGet(); }
public int getVues() { return vues.get(); }
\end{lstlisting}
\texttt{incrementAndGet()} est \textbf{atomique} : on obtient
\texttt{Vues : 200000} à chaque exécution, sans aucun \texttt{synchronized} (voir capture).

\capture{ex5b.png}{Compteur de vues avec \texttt{AtomicInteger} : 200000 (exact)}

\subsection{Question 5.5 — Tableau récapitulatif}
\begin{table}[H]\centering
\begin{tabular}{lcc}
\toprule
\textbf{Outil} & \textbf{Garantit la visibilité} & \textbf{Rend \texttt{c++} sûr} \\
\midrule
Rien & Non & Non \\
\texttt{volatile} & Oui & Non \\
\texttt{synchronized} & Oui & Oui \\
\texttt{AtomicInteger} & Oui & Oui \\
\bottomrule
\end{tabular}
\end{table}

% ================= EXERCICE 6 =================
\section{ExecutorService, Callable et Future}

\subsection{Question 6.1 — Un pool de 3 threads (\texttt{PoolSimple.java})}
Les 6 tâches sont exécutées dans un ordre variable par les threads du pool (voir
capture). \textbf{Combien de threads différents ?} \textbf{3}
(\texttt{pool-1-thread-1}, \texttt{-2}, \texttt{-3}), réutilisés pour plusieurs
tâches. \textbf{Si on oublie \texttt{pool.shutdown()} ?} Le programme \textbf{ne se
termine jamais} (threads du pool non daemon en attente).

\capture{ex6.png}{Un pool de 3 threads exécute 6 tâches}

\subsection{Question 6.2 — Callable et Future (\texttt{SommeFuture.java})}
\begin{lstlisting}
Future<Long> f = pool.submit(somme);
System.out.println("Résultat : " + f.get());
\end{lstlisting}
On obtient \texttt{Résultat : 500000500000} (\texttt{f.get()} est bloquant : il attend
la fin du calcul avant d'afficher, voir capture).

\capture{ex6b.png}{Résultat récupéré via \texttt{Future.get()} : 500000500000}

\subsection{Question 6.3 — invokeAll (\texttt{SommeParallele.java})}
Somme de 1 à 4\,000\,000 découpée en 4 morceaux.
\begin{lstlisting}
List<Future<Long>> resultats = pool.invokeAll(taches);
long total = 0;
for (Future<Long> f : resultats) {
    total += f.get();
}
\end{lstlisting}
On obtient \texttt{Total : 8000002000000} (voir capture).

\capture{ex6c.png}{Somme parallèle via \texttt{invokeAll} : 8000002000000}

\subsection{Question 6.4 — ScheduledExecutorService (\texttt{Bip.java})}
\begin{lstlisting}
ScheduledExecutorService pool = Executors.newScheduledThreadPool(1);
pool.scheduleAtFixedRate(() -> System.out.println("Bip"), 0, 1, TimeUnit.SECONDS);
Thread.sleep(5000);
pool.shutdown();
\end{lstlisting}
Affiche \texttt{Bip} toutes les secondes pendant 5\,s, puis s'arrête (voir capture).

\capture{ex6d.png}{« Bip » répété toutes les secondes pendant 5\,s}

% ================= EXERCICE 7 =================
\section{Le départ de la course}
5 coureurs (5 threads) s'échauffent, attendent un signal commun, puis partent
\textbf{en même temps}. L'arbitre (thread \texttt{main}) attend que tous soient prêts,
donne le départ, puis attend l'arrivée de tous. Classes : \texttt{LigneDeDepart}
(objet partagé), \texttt{Classement}, \texttt{Coureur} (thread), \texttt{Course}.

\subsection{Étape 1 — \texttt{LigneDeDepart} (TODO 1 à 4)}
\begin{lstlisting}
public synchronized void seDeclarerPret(String nom) {
    nbPrets++;
    System.out.println(nom + " est prêt (" + nbPrets + "/" + nbCoureurs + ")");
    notifyAll();                               // TODO 1
}
public synchronized void attendreTousPrets() throws InterruptedException {
    while (nbPrets < nbCoureurs) { wait(); }   // TODO 2
    System.out.println("Arbitre : tout le monde est prêt !");
}
public synchronized void donnerDepart() {
    departDonne = true;
    System.out.println(">>> PAN ! Départ ! <<<");
    notifyAll();                               // TODO 3
}
public synchronized void attendreDepart() throws InterruptedException {
    while (!departDonne) { wait(); }           // TODO 4
}
\end{lstlisting}

\subsection{Étapes 3 et 4 — \texttt{Coureur} et \texttt{Course}}
Points clés : \texttt{super(nom)} + initialisation des attributs ;
\texttt{seDeclarerPret} puis \texttt{attendreDepart} ; dans \texttt{Course}, boucle
\texttt{for} pour créer/démarrer les 5 coureurs, puis \texttt{join()} sur chacun.
La capture ci-dessous montre un déroulement complet.
\textbf{Vérifications :} aucun \texttt{s'élance} n'apparaît avant \texttt{PAN !} ; les 5
\texttt{s'élance} apparaissent juste après ; les positions vont de 1 à 5 sans doublon ;
\texttt{course terminée} s'affiche en dernier.

\capture{ex7.png}{Déroulement complet de la course}

\subsection{Étape 2 — Question sur \texttt{Classement}}
Sans \texttt{synchronized} dans \texttt{arriver()}, deux coureurs arrivant en même
temps pourraient lire la même \texttt{prochainePosition} → \textbf{doublon} et
position sautée.

\subsection{Étape 5 — Expériences}
\begin{enumerate}[nosep]
  \item \textbf{\texttt{notifyAll()} → \texttt{notify()} dans \texttt{donnerDepart()} :}
  \texttt{notify()} ne réveille qu'\textbf{un seul} coureur ; les autres restent
  bloqués, le \texttt{join()} ne se termine jamais → le programme ne se termine pas.
  \item \textbf{Supprimer \texttt{attendreTousPrets()} dans \texttt{Course} :}
  l'arbitre peut donner le départ avant que tous soient prêts. Les coureurs encore en
  échauffement \textbf{ne sont pas bloqués pour toujours} : en arrivant, leur
  \texttt{attendreDepart()} voit \texttt{departDonne == true} (grâce au \texttt{while})
  et ils partent aussitôt.
  \item \textbf{\texttt{while (!departDonne)} → \texttt{while (true)} :} la condition de
  sortie disparaît ; chaque coureur reste \textbf{bloqué pour toujours} dans
  \texttt{wait()} → la course ne démarre jamais.
\end{enumerate}

\subsection{Étape 6 — Comprendre les choix}
\begin{enumerate}[nosep]
  \item \textbf{Pourquoi \texttt{departDonne} est-il indispensable ?} C'est la
  \textbf{condition d'état} qui mémorise que le départ a eu lieu. Sans elle, un coureur
  qui fait \texttt{wait()} après un \texttt{notifyAll()} déjà passé attendrait un signal
  qui ne reviendra jamais (réveil manqué).
  \item \textbf{Pourquoi \texttt{while} et non \texttt{if} ?} À cause des réveils
  intempestifs (\emph{spurious wakeups}) et des réveils multiples : au réveil il faut
  \textbf{re-vérifier la condition}.
  \item \textbf{Un coureur dans \texttt{wait()} consomme-t-il du CPU ?} \textbf{Non} :
  \texttt{wait()} endort le thread et libère le verrou. Une attente active
  \texttt{while (!departDonne) \{\}} occuperait un cœur à 100 \% inutilement.
  \item \textbf{Pourquoi \texttt{notifyAll()} plutôt que \texttt{notify()} ?} Plusieurs
  threads attendent sur le \textbf{même objet} ; \texttt{notify()} n'en réveillerait
  qu'un (peut-être pas le bon), \texttt{notifyAll()} les réveille tous et chacun
  re-teste sa condition.
\end{enumerate}

\subsection{La version moderne — \texttt{CountDownLatch} (bonus)}
\texttt{LigneDeDepart} réécrite avec deux \texttt{CountDownLatch}, sans
\texttt{synchronized}, \texttt{wait()} ni \texttt{notifyAll()} :
\begin{lstlisting}
private final CountDownLatch prets;            // new CountDownLatch(5)
private final CountDownLatch depart = new CountDownLatch(1);
// coureur prêt            : prets.countDown();
// arbitre attend          : prets.await();
// arbitre donne le départ : depart.countDown();
// coureur attend le départ: depart.await();
\end{lstlisting}
Les classes \texttt{Coureur}, \texttt{Classement} et \texttt{Course} ne changent pas.
Même comportement, code beaucoup plus simple.

% ================= FAQ =================
\section{FAQ — Erreurs fréquentes rencontrées}
\begin{table}[H]\centering\small
\renewcommand{\arraystretch}{1.3}
\begin{tabular}{p{4.6cm}p{4.3cm}p{4.3cm}}
\toprule
\textbf{Symptôme} & \textbf{Cause} & \textbf{Solution} \\
\midrule
\texttt{unreported exception InterruptedException} & \texttt{sleep/join/wait/get} lèvent une exception vérifiée & \texttt{throws InterruptedException} ou \texttt{try/catch} \\
\texttt{IllegalThreadStateException} & \texttt{start()} appelé deux fois & créer un nouvel objet \texttt{Thread} \\
\texttt{IllegalMonitorStateException} & \texttt{wait()/notifyAll()} hors d'un bloc \texttt{synchronized} & appeler dans une méthode \texttt{synchronized} \\
variable non \texttt{final} dans une lambda & usage de \texttt{i} (variable de boucle) & copier avant : \texttt{int numero = i;} \\
\texttt{ArrayIndexOutOfBoundsException} (3.1) & c'est le bug à observer & le noter dans le tableau \\
Durée $\approx$ 0\,ms (ex. 1) & \texttt{join()} oublié & ajouter les \texttt{join()} \\
Le programme ne s'arrête jamais & pool non arrêté / \texttt{notifyAll()} oublié & \texttt{shutdown()} / \texttt{notifyAll()} \\
\bottomrule
\end{tabular}
\end{table}

% ================= CONCLUSION =================
\section*{Conclusion}
\addcontentsline{toc}{section}{Conclusion}
Ce TD a permis de manipuler les deux façons de créer des threads (\texttt{Runnable} et
\texttt{extends Thread}), d'observer leur cycle de vie, et surtout de comprendre les
\textbf{deux garanties distinctes} de la concurrence :
\begin{itemize}[nosep]
  \item la \textbf{visibilité} (\texttt{volatile}) — voir les écritures des autres threads ;
  \item l'\textbf{atomicité} (\texttt{synchronized}, \texttt{AtomicInteger}) — rendre
  indivisibles les opérations composées comme \texttt{c++} ou « vérifier puis ajouter ».
\end{itemize}
Les outils de haut niveau (\texttt{ExecutorService}, \texttt{Future},
\texttt{CountDownLatch}) simplifient grandement l'écriture de code concurrent correct.

\end{document}
